\section{USAMO 1973}
        \begin{enumerate}
        	\item (\textit{USAMO 1973}) Two points $P$ and $Q$ lie in the interior of a regular tetrahedron $ABCD$. Prove that angle $PAQ<60^o$.


 

	\item (\textit{USAMO 1973}) Let $\{X_n\}$ and $\{Y_n\}$ denote two sequences of integers defined as follows:


  $X_0=1$, $X_1=1$, $X_{n+1}=X_n+2X_{n-1}$ $(n=1,2,3,\dots),$

 $Y_0=1$, $Y_1=7$, $Y_{n+1}=2Y_n+3Y_{n-1}$ $(n=1,2,3,\dots)$.

 Thus, the first few terms of the sequences are:


  $X:1, 1, 3, 5, 11, 21, \dots$,

  $Y:1, 7, 17, 55, 161, 487, \dots$.

 Prove that, except for the "1", there is no term which occurs in both sequences.


 

	\item (\textit{USAMO 1973}) Three distinct vertices are chosen at random from the vertices of a given regular polygon of $(2n+1)$ sides. If all such choices are equally likely, what is the probability that the center of the given polygon lies in the interior of the triangle determined by the three chosen random points?


 

	\item (\textit{USAMO 1973}) Determine all the \href{https://artofproblemsolving.com/wiki/index.php?title=Root}{roots}, \href{https://artofproblemsolving.com/wiki/index.php?title=Real}{real} or \href{https://artofproblemsolving.com/wiki/index.php?title=Complex}{complex}, of the system of simultaneous \href{https://artofproblemsolving.com/wiki/index.php?title=Equation}{equations}


 $x+y+z=3$,
$x^2+y^2+z^2=3$,


 
$x^3+y^3+z^3=3$.

 

	\item (\textit{USAMO 1973}) Show that the cube roots of three distinct prime numbers cannot be three terms (not necessarily consecutive) of an arithmetic progression.


 

\end{enumerate}